\section{Performance Results}
\label{sec:results}

This section presents the results of a structured evaluation of \sys{} on
120 real user sessions collected over six weeks of platform deployment.
The sessions come from five academic disciplines and four language modes.
All numbers reported here are from real sessions — not synthetic benchmarks.

\subsection{Document Quality}

We measured document quality on a 0–100\% scale combining four factors:
whether the document compiles (35\% weight),
whether all citations are real and resolvable (35\% weight),
whether all figures in the text are present in the document (20\% weight),
and whether the document is written consistently in the specified language (10\% weight).

\begin{table}[h!]
\centering
\caption{Document quality by pipeline configuration, averaged across 120 sessions.}
\label{tab:quality}
\begin{tabular}{lrr}
\toprule
Configuration & Final quality score & Fabricated citation rate \\
\midrule
\sys{} full pipeline          & \textbf{97\%} & \textbf{0\%} \\
Without reference verification & 75\%          & 28\% \\
Without figures                & 88\%          & 0\% \\
Standard AI assistant (manual) & $\sim$60--70\% & 10--47\% \\
\bottomrule
\end{tabular}
\end{table}

The 97\% quality score means that the average \sys{} document:
compiles without any manual intervention,
has a fully verified bibliography,
has all figures present and correctly numbered, and
is written consistently in the specified language.
The remaining 3\% is primarily language consistency issues in Kazakh Latin sessions
(the newest and least common of the four supported modes).

\subsection{Quality by Academic Discipline}

\begin{table}[h!]
\centering
\caption{Final document quality by topic category.}
\label{tab:quality_by_topic}
\begin{tabular}{lrr}
\toprule
Topic category & Sessions & Quality score \\
\midrule
Computer Science \& Engineering    & 38 & 98\% \\
Education Research                  & 26 & 97\% \\
Biomedical Informatics              & 21 & 96\% \\
Social Sciences \& Economics        & 19 & 97\% \\
Mechanical / Civil Engineering      & 16 & 96\% \\
\midrule
\textbf{Overall}                   & \textbf{120} & \textbf{97\%} \\
\bottomrule
\end{tabular}
\end{table}

The consistency across disciplines — all five categories within 2 percentage points
of each other — is important.
It means \sys{} is not calibrated for one type of research at the expense of others.
Biomedical and engineering sessions score slightly lower because:
biomedical literature relies heavily on preprints (which are held to a slightly
different verification standard), and engineering documents require more complex
figures (detailed schematics) that are harder for the figure generator to compile.

\subsection{Figure Compilation Rate}

Every figure in a \sys{} document goes through an automatic compile-and-repair cycle.
The system attempts to generate a compilable figure, checks it against a \LaTeX{}
compiler, and if it fails, tries again up to 4 times total.

\begin{table}[h!]
\centering
\caption{Figure compilation success rate by attempt number.
         Based on 480 figures across 120 sessions.}
\label{tab:figures}
\begin{tabular}{lrr}
\toprule
After attempt \# & Figures compiled successfully & Cumulative success rate \\
\midrule
1 (first try)  & 335 of 480 & 69.8\% \\
2              & 427 of 480 & 89.0\% \\
3              & 470 of 480 & 97.9\% \\
4 (final)      & 477 of 480 & 99.4\% \\
\bottomrule
\end{tabular}
\end{table}

Only 3 of 480 figures (0.6\%) failed all 4 attempts.
All 3 were complex engineering circuit diagrams requiring a specialised library
not available in the current compiler environment — a known limitation that is
on the roadmap to resolve.

\subsection{Speed: Parallel Processing}

\sys{} collects references in parallel — multiple sources are fetched and verified
at the same time rather than one by one.
This is why Stage 2 (source collection and verification) takes minutes, not hours.

\begin{table}[h!]
\centering
\caption{Speed improvement from parallel reference fetching (Stage 2).}
\label{tab:speed}
\begin{tabular}{lrr}
\toprule
Number of parallel workers & Time for 20 references & Speed vs.\ sequential \\
\midrule
1 (sequential)  & 18.4 min & 1.0$\times$ \\
2               & 11.4 min & 1.6$\times$ \\
4 (default)     & 6.9 min  & 2.7$\times$ \\
6               & 5.6 min  & 3.3$\times$ \\
8               & 4.9 min  & 3.7$\times$ \\
\bottomrule
\end{tabular}
\end{table}

The platform uses 4 parallel workers by default, delivering 2.7$\times$ faster
reference collection than sequential processing.
More workers provide diminishing returns because some coordination overhead is
unavoidable regardless of parallelism.

\subsection{Cost Predictability}

Across all 120 sessions:

\begin{itemize}
  \item \textbf{Median session cost:} 30,920 platform micro-units (PMU)
        (approximately \$3.00--\$5.00 in API costs at current pricing)
  \item \textbf{95th percentile cost:} 61,200 PMU
  \item \textbf{Maximum cost:} 89,400 PMU — for a session with a budget ceiling
        of 120,000 PMU (well within budget)
  \item \textbf{Budget overruns:} \textbf{zero}. No session exceeded its pre-set budget ceiling.
\end{itemize}

The cost split between text generation (72\%) and figure generation (28\%)
shows that figures are the single most cost-effective addition: 4 figures
add 28\% to the total cost but contribute substantially to the final quality score.

\begin{figure}[h!]
\centering
\begin{tikzpicture}[
  kpi/.style={draw=#1, fill=#2, rounded corners=6pt,
              minimum width=3.6cm, minimum height=2.2cm,
              align=center, line width=1.5pt, inner sep=8pt},
  num/.style={font=\huge\bfseries, text=#1},
  lbl/.style={font=\scriptsize, text=MG, align=center, text width=3.2cm},
]
\node[kpi={C3}{L3}] (K1) at (-6.5, 0) {};
\node[num={C3}]     at (-6.5, 0.35) {97\%};
\node[lbl]          at (-6.5,-0.55) {mean document\\quality score};

\node[kpi={C1}{L1}] (K2) at (-2.2, 0) {};
\node[num={C1}]     at (-2.2, 0.35) {0\%};
\node[lbl]          at (-2.2,-0.55) {fabricated citation\\rate};

\node[kpi={C2}{L2}] (K3) at ( 2.2, 0) {};
\node[num={C2}]     at ( 2.2, 0.35) {99.4\%};
\node[lbl]          at ( 2.2,-0.55) {figure compilation\\success (4 tries)};

\node[kpi={C5}{L5}] (K4) at ( 6.5, 0) {};
\node[num={C5}]     at ( 6.5, 0.35) {2.7×};
\node[lbl]          at ( 6.5,-0.55) {speed gain from\\parallel fetching};

\node[kpi={C6}{L6}] (K5) at (-4.4,-3.0) {};
\node[num={C6}]     at (-4.4,-2.65) {100\%};
\node[lbl]          at (-4.4,-3.55) {sessions within\\budget ceiling};

\node[kpi={C4}{L4}] (K6) at ( 0.0,-3.0) {};
\node[num={C4}]     at ( 0.0,-2.65) {20--50};
\node[lbl]          at ( 0.0,-3.55) {minutes per complete\\document};

\node[kpi={C7}{L7}] (K7) at ( 4.4,-3.0) {};
\node[num={C7}]     at ( 4.4,-2.65) {120};
\node[lbl]          at ( 4.4,-3.55) {real sessions in\\evaluation corpus};

\node[font=\small\itshape, text=MG] at (0,-4.8)
  {All metrics measured on real user sessions across 5 disciplines and 4 languages.};
\end{tikzpicture}
\caption{Key performance indicators from the \sys{} platform evaluation.
         All seven metrics come from 120 real user sessions collected over
         six weeks of deployment — not from synthetic benchmarks.}
\label{fig:results_overview}
\end{figure}

